\documentclass[11pt]{article}
\usepackage[margin=1in]{geometry}
\usepackage[hidelinks]{hyperref}
\usepackage{enumitem}

\title{PA1 -- Perceptron}
\date{Due: April 23, 2026 at 11:59 PM}

\begin{document}
\maketitle

\section*{Overview}
In this assignment you will implement binary and multiclass perceptrons, apply
simple feature engineering, and evaluate performance on the handwritten digits
(8x8) dataset. All graded tasks are marked with TODO comments in
\texttt{PA1.ipynb}.

\section*{Learning goals}
\begin{itemize}[leftmargin=1.25em]
  \item Implement the perceptron update rule for binary and multiclass settings.
  \item Practice basic data preprocessing and evaluation metrics.
  \item See how feature mapping can make XOR linearly separable.
  \item Interpret training behavior using a learning curve.
\end{itemize}

\section*{Rules and allowed libraries}
\begin{itemize}[leftmargin=1.25em]
  \item Use only Python standard libraries for training logic.
  \item It is OK to use sklearn to load the dataset.
  \item GenAI tools that generate output may be used in course assignments (not tests or exams) to support your learning, as long as you do so honestly through proper documentation, citation, and acknowledgement.
\end{itemize}

\section*{Deliverables}
\begin{itemize}[leftmargin=1.25em]
  \item Submit \texttt{PA1.ipynb} with all TODOs completed.
  \item Submit the notebook to Gradescope as an \texttt{.ipynb} file, not as a PDF.
  \item Do not rename required functions or change function signatures.
\end{itemize}

\section*{Point values (100 total)}

\subsection*{Question 0: Setup + Utils (10 points)}
\begin{itemize}[leftmargin=1.25em]
  \item 0.1 Load the dataset (\texttt{load\_digits\_data}) -- 5 pts
  \item 0.2 Normalize features (\texttt{normalize\_features}) -- 5 pts
\end{itemize}

\subsection*{Question 1: Binary Perceptron (35 points)}
\begin{itemize}[leftmargin=1.25em]
  \item 1.1 Initialize parameters (\texttt{BinaryPerceptron.create}) -- 5 pts
  \item 1.2 Score an example (\texttt{BinaryPerceptron.score}) -- 5 pts
  \item 1.3 Predict a label (\texttt{BinaryPerceptron.predict}) -- 5 pts
  \item 1.4 Update on a mistake (\texttt{BinaryPerceptron.update}) -- 10 pts
  \item 1.5 Train over epochs (\texttt{BinaryPerceptron.fit}) -- 10 pts
\end{itemize}

\subsection*{Question 2: Binary Classification Task (20 points)}
\begin{itemize}[leftmargin=1.25em]
  \item 2.1 Filter to two digits (\texttt{make\_binary\_dataset}) -- 8 pts
  \item 2.2 Binary accuracy (\texttt{binary\_accuracy}) -- 6 pts
  \item 2.3a XOR dataset (\texttt{xor\_dataset}) -- 3 pts
  \item 2.3b XOR feature map (\texttt{map\_xor\_features}) -- 3 pts
\end{itemize}

\subsection*{Question 3: Multiclass Perceptron (25 points)}
\begin{itemize}[leftmargin=1.25em]
  \item 3.1 Initialize parameters (\texttt{MulticlassPerceptron.create}) -- 5 pts
  \item 3.2 Compute scores (\texttt{MulticlassPerceptron.scores}) -- 5 pts
  \item 3.3 Predict class (\texttt{MulticlassPerceptron.predict}) -- 5 pts
  \item 3.4 Update on a mistake (\texttt{MulticlassPerceptron.update}) -- 5 pts
  \item 3.5 Train over epochs (\texttt{MulticlassPerceptron.fit}) -- 5 pts
\end{itemize}

\subsection*{Question 4: Evaluation + Learning Curve (5 points)}
\begin{itemize}[leftmargin=1.25em]
  \item 4.1 Multiclass accuracy (\texttt{multiclass\_accuracy}) -- 2 pts
  \item 4.2 Mistakes per epoch (\texttt{train\_binary\_with\_mistakes}) -- 3 pts
  \item 4.3 Plotting (provided) -- 0 pts
\end{itemize}

\subsection*{Question 5: Short Written Responses (5 points, manually graded)}
\begin{itemize}[leftmargin=1.25em]
  \item Written responses in the notebook -- 5 pts
\end{itemize}

\section*{Notes}
\begin{itemize}[leftmargin=1.25em]
  \item Train/test split (provided utility) is provided.
  \item The autograded portion of PA1 is worth 95 points total.
  \item Gradescope will show some autograder checks immediately after submission, and additional autograder checks also run during grading but are not shown in advance.
  \item Question 5 is graded manually from the written responses in the notebook.
\end{itemize}

\section*{Submission checklist}
\begin{itemize}[leftmargin=1.25em]
  \item All TODOs are implemented.
  \item Notebook runs end-to-end without errors.
  \item You did not rename required functions or change function signatures.
  \item Your \texttt{.ipynb} file is uploaded to Gradescope and the autograder shows no submission errors.
\end{itemize}

\section*{Tips}
\begin{itemize}[leftmargin=1.25em]
  \item Start with the small utility functions (Question 0), then build the binary
  perceptron (Question 1) before moving to multiclass.
  \item For shuffling, keep X and y aligned by shuffling indices.
  \item Use small digit pairs (e.g., 0 vs 1) to sanity-check your binary model.
\end{itemize}

\end{document}
